% Generated draft: requires mathematical and source review.
\documentclass[11pt]{article}
\usepackage[margin=25mm]{geometry}
\usepackage{fontspec}
\setmainfont{Latin Modern Roman}
\usepackage{amsmath,amssymb,mathtools,unicode-math}
\setmathfont{Latin Modern Math}
\usepackage{enumitem,needspace,xurl,hyperref}
\setlength{\parindent}{0pt}
\setlength{\parskip}{5pt}
\setlength{\emergencystretch}{3em}
\newcommand{\N}{\mathbb{N}}
\newcommand{\Z}{\mathbb{Z}}
\newcommand{\Q}{\mathbb{Q}}
\newcommand{\R}{\mathbb{R}}
\newcommand{\C}{\mathbb{C}}
\newcommand{\F}{\mathbb{F}}
\newcommand{\A}{\mathbb{A}}
\newcommand{\PP}{\mathbb{P}}
\newcommand{\E}{\mathbb{E}}
\newcommand{\eps}{\varepsilon}
\newcommand{\dd}{\,\mathrm{d}}
\begin{document}
\section*{UnsolvedMath ID: 11000080}
\textbf{Source problem code:} catalogue entry
\subsection{Definitions and mathematical statement}
For a pseudo-Anosov mapping class \(f\in\operatorname{Mod}_g\), let \(\lambda(f)>1\) be its dilatation, the expansion factor for its invariant measured foliations. Write
\[
 \delta_g=\inf\{\lambda(f):f\in\operatorname{Mod}_g
                         \text{ pseudo-Anosov}\}.
\]
Question 7.4 asks whether the classes attaining this minimum form one conjugacy class:
\[
 \left|\{[f]_{\rm conj}: f\text{ pseudo-Anosov},\ \lambda(f)=\delta_g\}\right|=1\ ?
\]
Conjugacy is in the source's \(\operatorname{Mod}_g\), not an unspecified larger equivalence allowing inversion, central multiplication or orientation reversal. The genus range and exact surface conventions are not supplied.

The source claims uniqueness for \(g=1\), represented by
\[
 A=\begin{pmatrix}2&1\\1&1\end{pmatrix}.
\]
This needs qualification under the standard \(\operatorname{Mod}_1\cong\mathrm{SL}_2(\Z)\) convention: \(A\) and \(-A\) have the same expansion-factor modulus but different traces, so are not conjugate in that group. Thus the stated genus-one example may be using a sign quotient or an extra foliation/action convention. It is retained as a source claim rather than used as an unqualified base case.

The subsequent context proposes refining the minimum-dilatation question by singularity data. Fix a genus and an admissible tuple \((k_1,\ldots,k_r)\) recording the prong numbers at the stable foliation's \(r\) singularities. In the usual closed genus-at-least-two convention, singularities have \(k_i\ge3\) and Poincaré--Hopf gives \(\sum_i(k_i-2)=4g-4\), hence finitely many possible tuples. These necessary constraints are not asserted sufficient for the source's notion of admissibility. The source cites Masur--Smillie [MS] for realization of every admissible tuple by some pseudo-Anosov stable foliation. It refers to \(L(\operatorname{Mod}_g)\) without defining whether that means dilatation or its logarithm, and then introduces a following unprovided question. No missing refinement statement is invented.
\subsection{Short English statement}
Is the least-dilatation pseudo-Anosov unique up to the stated conjugacy relation? The source's genus-one uniqueness claim needs an additional sign or equivalence convention. The following admissible-prong-data discussion is retained as a refinement program, without replacing the main uniqueness question.
\subsection{Sources}
\begin{itemize}
\item Local UnsolvedMath registry, permanent ID \(11000080\), exact source code \texttt{catalogue entry}, Question 7.4, PDF page 57. No individual proposer or proposed year is supplied.
\item Farb, edited volume \emph{Problems on Mapping Class Groups and Related Topics}: \url{https://www.math.uchicago.edu/~farb/papers/mcgbook.pdf}. The volume has 329 explicitly typeset numbered items, repeated local numbering distinguished by PDF page. The extract ends with chapter/page furniture and an introduction to another question.
\item Internal reference Masur--Smillie [MS] for stable-foliation singularity-data realization; the full entry and admissibility hypotheses are absent.
\item Extraction pdf-font-structure, accessed 2026-07-29. Source evidence NEEDS\_REVIEW presents an open problem and requires a current primary-literature audit. Supplied OPEN-TRIAGE summary says some small-genus cases may have uniqueness and general uniqueness is open, without exact cases or references.
\item Dataset attribution: Ulam AI / UnsolvedMath Contributors, CC BY 4.0. Source rights NEEDS\_REVIEW: publicly accessible author/publisher copy, redistribution terms requiring review.
\end{itemize}
\end{document}
